\documentclass[11pt,letterpaper]{article}
\usepackage[T1]{fontenc}
\usepackage{lmodern,amsmath,amssymb,amsthm,mathtools}
\usepackage[margin=0.9in]{geometry}
\usepackage{booktabs,longtable,array,microtype,xcolor,enumitem,fancyhdr}
\usepackage[hidelinks]{hyperref}
\definecolor{ink}{HTML}{18384C}
\definecolor{muted}{HTML}{53636D}
\setlength{\parindent}{0pt}
\setlength{\parskip}{0.48em}
\setlength{\emergencystretch}{2em}
\setlist{nosep,leftmargin=1.6em}
\pagestyle{fancy}
\fancyhf{}
\fancyhead[L]{\small\color{muted}G(7) investigation}
\fancyhead[R]{\small\color{muted}22 November 2026}
\fancyfoot[C]{\small\thepage}
\setlength{\headheight}{14pt}
\newtheorem{proposition}{Proposition}[section]
\newtheorem{lemma}[proposition]{Lemma}
\newtheorem{theorem}[proposition]{Theorem}
\theoremstyle{remark}
\newtheorem*{scope}{Scope}
\newcommand{\A}{\mathcal A}
\newcommand{\M}{\mathfrak M}
\newcommand{\m}{\mathfrak m}
\newcommand{\HH}{\mathcal H}
\newcommand{\eps}{\varepsilon}
\newcommand{\dd}{\,d\alpha}
\newcommand{\e}{\mathrm e}
\newcommand{\R}{\mathbb R}
\newcommand{\Z}{\mathbb Z}
\newcommand{\ReportState}{Completed research record}
\newcommand{\RunText}{The authorized research run lasted 3 hours, 0 minutes and 17 seconds, from 11:11:18 AM EST to 2:11:35 PM EST on 22 November 2026.}
\newcommand{\BoundGSeven}{31}
\newcommand{\CurrentGap}{0.015452339536143}
\newcommand{\CurrentGapExact}{\frac{108166376753}{7000000000000}}
\newcommand{\GapReductionPercent}{41.53}
\newcommand{\CurrentDeltaTwentyEight}{0.093564057962}
\newcommand{\CurrentDeltaTwentyTwo}{0.423201151219}
\newcommand{\CurrentDeltaTwenty}{0.605567600000}
\newcommand{\CurrentDeltaThirty}{0.015452339536143}
\newcommand{\CurrentDeltaThirtyExact}{\frac{108166376753}{7000000000000}}
\newcommand{\CurrentDeltaThirtyTwoStar}{-0.025814827130524}
\newcommand{\CurrentDeltaThirtyTwoStarExact}{\frac{-542111369741}{21000000000000}}
\newcommand{\CurrentRealZero}{30.748892680758}
\newcommand{\CurrentRealZeroExact}{\frac{13323656630259}{433305250000}}
\newcommand{\CoefficientThree}{3.011012515971}
\newcommand{\CoefficientFour}{4.084400784447}
\newcommand{\CoefficientFive}{5.219004846379}
\newcommand{\ConeGapShort}{0.0688099376723}
\newcommand{\ConeGapShortNine}{0.0768987937077}
\newcommand{\ConeGapAll}{0.015452339536143}
\newcommand{\ConeGapAllNine}{0.0191363819577}
\newcommand{\CertificateName}{\path{moment_refinement/quadratic_height_distribution_exact_certificate.json}}
\newcommand{\CertificateRounds}{220}
\newcommand{\GeneratedAt}{2026-11-22 19:11:46 UTC}
\newcommand{\MomentTableRows}{22 & 0.4242973 & 0.423201151219 \\
24 & 0.2891303 & 0.283697712880 \\
26 & 0.1843399 & 0.176197101178 \\
28 & 0.1053935 & 0.093564057962 \\
30 & 0.0484598 & 0.015452339537 \\
32 & 0.0051012 & 0.000000000000 \\}
\newcommand{\ParameterRows}{22 & $[10,9,8,7,6]$ & \shortstack[l]{0.127866478, 0.120048112 \\ 0.108136229, 0.084376799 \\ 0.026837914} & \shortstack[l]{Quadratic sampling;\\ auxiliary mean} \\
24 & $[11,10,10,10]$ & \shortstack[l]{0.130838206, 0.124245584 \\ 0.115637814, 0.096791920} & Quadratic sampling \\
26 & $[12,11,10,10]$ & \shortstack[l]{0.132885530, 0.126603736 \\ 0.115836987, 0.097227991} & Quadratic sampling \\
28 & $[13,12,12]$ & \shortstack[l]{0.134436321, 0.128133606 \\ 0.118200529} & Quadratic sampling \\
30 & $[14,13,12]$ & \shortstack[l]{0.135787759, 0.129473310 \\ 0.118256158} & \shortstack[l]{Available policy;\\ final bound recycled} \\}

\begin{document}
\thispagestyle{empty}
{\color{muted}\small RESEARCH CONTINUATION\hfill\ReportState}\par
\vspace{0.7cm}
{\color{ink}\LARGE\bfseries G(7): a three-hour investigation}\par
\vspace{0.15cm}
{\large Exact auxiliary progress and the remaining record barrier}\par
\vspace{0.35cm}
\textbf{Record status.} The verified unconditional bound remains
\[
                         \boxed{G(7)\le\BoundGSeven.}
\]
This investigation does not establish $G(7)\le30$. It improves and certifies several
auxiliary estimates, identifies precise sufficient inputs for a record improvement,
and records where the remaining arguments stop.

\RunText

The best recorded deficit in the tested thirty-variable argument is
\[
 D=\CurrentGapExact=\CurrentGap\ldots>0.
\]
The extension began with deficit $0.0264290504625$. The reduction is
\GapReductionPercent\%, within this particular analytic calculation. It is not a
probability of success or a universal measure of the difficulty of Waring's problem.
The initial best test used seven ordinary and 23 smooth variables; the current
minimum is attained with four ordinary and 26 smooth variables. The percentage
compares the best tested deficits at those two stages.

\textbf{Deliverables.} This report is accompanied by its TeX source, exact rational
certificates, the scripts that verify them, independent analytic audits, and a research
archive. Numerical searches choose parameters; the certified conclusions use finite
rational evaluations and explicit analytic inequalities.

\textbf{Scope.} Only unconditional classical $G(7)$ is investigated. Work on $G(4)$,
GRH or ERH variants, Waring--Siegel variants, and Fekete questions is outside this run.
No priority or publication novelty claim is made for the auxiliary deductions.

\vfill
{\small\color{muted}Report data generated: \GeneratedAt.\par
Current exact seventh-power certificate:\par
\CertificateName.}
\newpage
\tableofcontents
\newpage

\section{Meaning of the calculation}

For a positive integer $k$, $G(k)$ is the least $s$ such that every sufficiently
large integer is a sum of at most $s$ positive integral $k$th powers. Wooley's established result
$G(7)\le31$ is the record used as the stopping threshold in this run \cite{W16}.
An early stop would require a complete checked proof of $G(7)\le30$ or better.
An improved mean value, a negative floating-point objective, or a conditional
representation theorem alone would not meet that condition.

Write
\[
 f_P(\alpha)=\sum_{1\le x\le P}\e(\alpha x^7),\qquad
 g_Q(\alpha)=\sum_{x\in\A(Q,R)}\e(\alpha x^7),\qquad
 \e(t)=\exp(2\pi i t).
\]
Here $\A(Q,R)$ consists of the positive integers at most $Q$ whose prime factors are
at most $R$. An admissible excess $\Delta_v$ means that, for every $\eps>0$ and a
sufficiently small positive smoothness exponent,
\begin{equation}\label{eq:admissible}
 \int_0^1|g_Q(\alpha)|^v\dd\ll Q^{v-7+\Delta_v+\eps}.
\end{equation}
For even order $v=2s$ the traditional notation is
$\lambda_s=2s-7+\Delta_{2s}$. Degree-fourteen coefficient moments use the analogous
notation $\lambda_a^{(14)}=2a-14+\Delta_{2a}^{(14)}$.

All uses of $P^\eps$ and small powers of $R$ follow the finite-policy convention:
first fix the finitely many parameters and inequalities, then choose the smoothness
exponent small enough for all of them. This is distinct from discarding a fixed
positive power of $P$.

\subsection{Current certified smooth moments}

The baseline column is the recovered checkpoint, rounded upward as recorded there.
The current column is also rounded upward; exact rational values are in the certificate.
\begin{center}
\begin{tabular}{rrr}
\toprule
Order $v$ & Recovered excess & Current admissible excess\\
\midrule
\MomentTableRows
\bottomrule
\end{tabular}
\end{center}
The order-$32$ conclusion can be strengthened to a bound without an $\eps$ loss;
Section~\ref{sec:recycle} gives the argument. The order-$30$ excess is
$\CurrentDeltaThirtyExact$, or smaller if a later exact certificate improves it.

\subsection{Why the remaining sign matters}

The established ordinary minor-arc mean is
\begin{equation}\label{eq:f56}
 \int_{\m_P(P)}|f_P(\alpha)|^{56}\dd\ll P^{48+\eps}.
\end{equation}
This is the degree-seven case of the modern Vinogradov mean-value input used in
\cite[Lemma~3.1]{W16}. For four ordinary and 26 smooth variables, H\"older gives
\begin{align}
 \int_{\m_P(P)}|f_P|^4|g_P|^{26}\dd
 &\le\left(\int_{\m_P(P)}|f_P|^{56}\dd\right)^{1/14}
       \left(\int_0^1|g_P|^{28}\dd\right)^{13/14}\notag\\
 &\ll P^{23+(13\Delta_{28}-1)/14+\eps}.\label{eq:gap}
\end{align}
The expected representation main term has size $P^{23}$. Thus a positive value
of $(13\Delta_{28}-1)/14$ does not make the minor arcs smaller than the main term.
In this allocation a sufficient improvement would be $\Delta_{28}<1/13$.
The present $\Delta_{28}=\CurrentDeltaTwentyEight$ does not satisfy that inequality.

\section{The coefficient large-sieve refinement}

The auxiliary sums $F_j$ arise from repeated efficient differencing
\cite[Section~2]{VW00}. Put $d=7-j$, and write
\[
 M_i=P^{\phi_i},\quad H_i=P M_i^{-7},\quad T=\sum_{i=1}^j\phi_i,
 \quad Q=P^{1-T},\quad A=P\widetilde H_j\widetilde M_j.
\]
Here $0\le\phi_i\le1/7$, $\widetilde H_j=\prod H_i$, and
$\widetilde M_j=\prod M_i$.
The centered differenced polynomial has leading terms proportional to
\[
 (\textstyle\prod h_i)
 \left(z^d+c z^{d-2}\sum_i h_i^2m_i^{14}+\cdots\right).
\]
After the prescribed ordinary differences, a coefficient block has the form
\[
 E_a(\alpha)=\sum_{h\le H}\sum_{z\le J}
       \left|\sum_{m\in C(M)}\e(C\alpha z h^3m^{14})\right|^{2a},
 \qquad C(M)=\A(MR,R)\cap(M,MR],
\]
where $C$ is a fixed nonzero integer. Up to fixed factors and powers of $R$,
$JH^3M^{14}=Q^7$.

\subsection{Stopping the original proof at its internal approximant}

The proof of Vaughan--Wooley's large-sieve lemma first constructs its own reduced
approximant $b/t$. With $N$ a support bound and $Q_0=H^3JN$, put
\[
 L=t+Q_0|t\alpha-b|.
\]
Before the final matching to an externally prescribed approximant, the admissible
cutoff is
\[
 Y\ll\min\{N,J,(JN/H^3)^{1/2}\}.
\]
The original additional conditions involving $Y^3$ and the external denominator
cutoff are used in that final matching step. They are unnecessary for a conclusion
expressed through this internal $L$, or through a lower bound valid for every
rational approximant. This stopping point was checked directly against the full
proof of \cite[Lemma~4.5]{VW95} by three readers in this run.

For a divisor $t_0\mid t$, define $\ell$ to be the least positive integer with
$t_0\mid\ell^3$. A nonempty class has $\ell\le H$ and $t_0\le\ell^3$.
Writing $\beta=\alpha-b/t$, its reciprocal contribution is at most
\begin{equation}\label{eq:divisorclass}
 \frac{t_0}{t}\sum_{z\le H/\ell}
              (1+JN\ell^3z^3|\beta|)^{-1}.
\end{equation}
Trivial summation bounds this by $H^3/t$. When $\beta\ne0$, comparison with
$\sum z^{-3}$ bounds it by $C/(tJN|\beta|)$. Combining the two bounds gives
\[
 \eqref{eq:divisorclass}\ll \frac{H^3}{t+Q_0|t\alpha-b|}.
\]
The first bound alone handles $\beta=0$. Summing the divisor classes costs
$P^\eps$. Retaining the original cube-root estimate as well gives
\begin{equation}\label{eq:reciprocal}
 \text{reciprocal contribution}\ll
 P^\eps\min\{H L^{-1/3},H^3/L\}.
\end{equation}
The separate $H/Y$ term of the large-sieve argument remains present. No assumption
$t\le H$, $t_0\le H$, or $q\le H^3$ has been introduced.

\subsection{A useful direct version}

Applying the original large-sieve lemma separately at each $h$, without the averaging
corollary, also proves on the auxiliary minor arcs
\[
 E_a\ll P^\eps H(Q^7/P+J+M^{14})
              M^{2a-14+\Delta_{2a}^{(14)}}.
\]
Relative to $JHM^{2a}$, the saving is
\[
 \min(21\phi-2,14\phi,4-7T+7\phi)
                         -\Delta_{2a}^{(14)}\phi.
\]
This was the first improvement in the extension. The height-dependent refinement
below subsequently gives the stronger current certificate.

\subsection{Direct grouping of the cubic product}

A further sampling estimate avoids the square-root cutoff in the preceding
averaging argument. Let $g$ have frequency support in an interval of length $N$,
let $p\ge1$ be fixed, and put $V=JH^3N$. For a reduced rational approximant write
$L=q+V|q\alpha-a|$. The reciprocal-sum estimate and its elementary transference
form give
\begin{equation}\label{eq:grouped-sampling}
 \sum_{h\le H}\sum_{z\le J}|g(\alpha h^3z)|^p
 \ll P^\eps\left(HJ+\frac{V}{L}+H^3N+L\right)
                      \int_0^1|g(\beta)|^p\,d\beta.
\end{equation}
For an intrinsic minimizing approximant, $L\ll\sqrt V$, so the final term is
absorbed into $V/L$.

Here is a proof that also covers real $p$. First shift the frequency interval,
which leaves $|g|$ unchanged. The de la Vall\'ee Poussin kernel
$V_N=2K_{2N}-K_N$ reproduces $g$, has bounded $L^1$ norm, and is dominated in
absolute value by the positive kernel $2K_{2N}+K_N$. Jensen's inequality reduces
the assertion to a uniform bound for the sum of translates of that kernel.
Fourier expansion followed by the sum in $z$ bounds it by
\[
 HJ+\sum_{h\le H}\sum_{1\le b\le2N}
                   \min\{J,\|\alpha h^3b\|^{-1}\}.
\]
Group the product $n=h^3b$. Its multiplicity is at most the divisor function,
hence $P^\eps$. The elementary reciprocal-sum estimate behind
\cite[Lemma~4.4]{VW95}, at length $O(H^3N)$, gives
$P^\eps(V/L+H^3N+J+L)$; absorb $J$ into $HJ$.
This proves \eqref{eq:grouped-sampling} without expanding a fractional power
or assuming that its Fourier coefficients are nonnegative.

With $N\asymp(M_iR)^{14}$ and
$J=P^{5-j}\prod_{l\ne i}H_l$, normalization by $H_iJM_i^{2a}$ gives
the additional saving
\begin{equation}\label{eq:grouped}
 b_{i,a}^{\rm group}(\ell)
 =\min\{14\phi_i,\,2-7T+21\phi_i,\,\ell-2+14\phi_i\}
                         -\Delta_{2a}^{(14)}\phi_i.
\end{equation}
The new middle cap exceeds the old square-root cap by $(3-7T)/2$.
The other caps still matter, so this is retained as a separate option alongside
the earlier bounds. Fixed multipliers and the powers of $R$ in the actual
bandwidth are handled by the same finite-policy smoothness convention.

\section{Coupling both estimates to the same rational height}

Define the intrinsic height
\begin{equation}\label{eq:height}
 \HH_Q(\alpha)=\min_{(a,q)=1}\max\{q,Q^7|q\alpha-a|\}.
\end{equation}
Its sublevel set is the standard major-arc union $\M_Q(Z)$, and Dirichlet gives
$\HH_Q\ll Q^{7/2}$. The original efficient-differencing major arcs include powers
of $R$ in their widths. The needed fact is the inclusion
$\M_Q(cPR^{-B})\subseteq\M_{\rm auxiliary}$ for a fixed sufficiently large $B$.
Thus the auxiliary minor region has height at least $P^{1-O(\eta)}$.

The quadratic-difference factor retains its old bound
$D_j\ll P^\eps Q^7/P$ when $Q^7>P^2$. Its plateau is not replaced by an
unsupported extra height saving. For coefficient block $i$, put
\[
 h_i^{\rm exp}=1-7\phi_i,\qquad
 y_i=\min\left\{14\phi_i,4-7T+7\phi_i,
                         \frac{1-7T+42\phi_i}{2}\right\}.
\]
On a shell $\HH_Q\asymp P^\ell$, two high-moment savings, in addition to
\eqref{eq:grouped}, are
\begin{align}
 b_{i,a}^{\rm ref}(\ell)
   &=\min\{y_i,\ell-2h_i^{\rm exp}\}
                              -\Delta_{2a}^{(14)}\phi_i,\label{eq:ref}\\
 b_{i,a}^{\rm cube}(\ell)
   &=\min\{y_i,\ell/3\}
                              -\Delta_{2a}^{(14)}\phi_i.\label{eq:cube}
\end{align}
The first-moment coefficient bounds remain available. For $j\le3$ their saving is
$\phi_i$; for $j=4,5$ it is
\[
 \min\left\{\phi_i,1-7T+28\phi_i,(1-T)/4,
                          [7(1-T)-27\phi_i]/6\right\}.
\]

Choose nonnegative H\"older weights $w_{i,a}$ with
\[
 \sum_{i,a}w_{i,a}\le1,\qquad
 \sum_a a w_{i,a}\le1\quad\text{for each }i.
\]
Missing powers are bounded by the coefficient-sum lengths, and unused total weight
is assigned to a constant factor. If $S(\ell)$ is the weighted saving, the auxiliary
pointwise estimate is
\begin{equation}\label{eq:pointwise}
 |F_j(\alpha)|\ll P^\eps A
     P^{-(1+\min\{1,S(\ell)\})/2^{d-1}}.
\end{equation}
The cap at1 explicitly retains the ordinary-differencing diagonal term.

The crucial complementary estimate is the restricted smooth mean
\cite[Theorem~4.2]{BW23}:
\begin{equation}\label{eq:restricted}
 \int_{\M_Q(Z)}|g_Q|^v\dd
             \ll Q^{v-7+\eps}Z^{2\Delta_v/7},\qquad v\ge2.
\end{equation}
Multiplying \eqref{eq:pointwise} and \eqref{eq:restricted}, then summing logarithmically
many shells, bounds the auxiliary minor integral by
\begin{equation}\label{eq:terminal}
 A Q^{v-7}P^{\eps+max_{1\le\ell\le7(1-T)/2}
 \{2\Delta_v\ell/7-(1+\min(1,S(\ell)))/2^{d-1}\}}.
\end{equation}
The lower endpoint differs from1 by $O(\eta)$, which is handled by the stated
smoothness convention. The final constant-factor top shell uses the full smooth
moment and has the same exponent. Fixed phase multipliers alter heights only by
bounded factors and harmless powers of $R$.

The old terminal combined the weakest coefficient estimate with the full smooth
moment. Formula~\eqref{eq:terminal} keeps their dependence on height together.
A region with a weaker coefficient saving also has a smaller available smooth mean.

\subsection{Why finitely many exact endpoint checks suffice}

For one fixed recipe using \eqref{eq:ref}, \eqref{eq:cube}, and
\eqref{eq:grouped} as separate options,
each saving is concave in $\ell$. Hence $\min(1,S(\ell))$ is concave and the exponent
in \eqref{eq:terminal} is convex. Its maximum on an interval occurs at an endpoint.
Different recipes may be frozen on finitely many closed intervals that cover the
whole height range. The certificate checks both endpoints of every interval.

This shortcut would be invalid for an unchecked optimized envelope that changes
recipes internally. The actual certificates specify the intervals and recipes,
including all rational H\"older weights, explicitly.

\subsection{Retaining the quadratic fourth moment}\label{sec:quadratic-fourth}

A further improvement keeps the quadratic and coefficient factors together until
after the ordinary differences. Put $r=5-j$ and
$B=\widetilde H_j\widetilde M_j$. The unseparated Weyl inequality is
\begin{equation}\label{eq:unseparated}
 |F_j|^{2^r}\ll P^{2^r-1}B^{2^r}
            +P^{2^r-r-1}B^{2^r-1}|G_r|.
\end{equation}
Here $G_r$ still sums over one original tuple of coefficient variables $\mathbf m$.
For a direct two-difference source example see \cite[(12.17)]{VW95}; the factor
separation occurs in its next equation. More generally, Cauchy over the existing
parameter tuple gives
$|G_u|^2\ll P^{2u+1}B^2+P^uB|G_{u+1}|$, which proves
\eqref{eq:unseparated} inductively. Precisely $d-2$ differences turn the original
$z^{d-2}\sum h_i^2m_i^{14}$ term into a constant that remains inside $G_r$.

The resulting quadratic phase is a fixed multiple of
\[
 \Gamma\left(3\xi^2+\sum_{i\le j}h_i^2m_i^{14}
                         +\sum_{\nu\le r}l_\nu^2\right),\qquad
 \Gamma=\prod_{i\le j}h_i\prod_{\nu\le r}l_\nu.
\]
The $l_\nu$ are ordinary shifts. Complete the shift-dependent $\xi$ interval and
its fixed parity condition by Fourier inversion, using one common frequency $t$.
The completion kernel has $L^1$ norm $O(\log P)$. It is therefore enough to bound
\[
 \sum_{\text{shifts}}|Z_t(C_0\Gamma\alpha)|
              \prod_{i\le j}|L_i(C_i\Gamma h_i^2\alpha)|,
 \quad Z_t(\beta)=\sum_{|z|\le CP}\e(\beta z^2+tz),
\]
uniformly in $t$. The $L_i$ are the full smooth-interval coefficient sums.
Taking absolute values of the individual coefficient sums before the ordinary
differences would lose the information used here.

The product scale of all shifts is
$H_*=P^{5-j}\widetilde H_j=P^{5-7T}$, and $P^2H_*=Q^7$.
Quadratic differencing, uniformly in $t$, gives on the auxiliary minor region
\[
 \sum_{\text{shifts}}|Z_t(C_0\Gamma\alpha)|^2
       \ll P^{1+\eps}H_*.
\]
For the fourth moment, $Z_t^2$ has bandwidth $O(P^2)$ and
$\int|Z_t|^4\ll P^{2+\eps}$ uniformly in $t$, by orthogonality and the
difference-of-squares count. Grouping $\Gamma$ with divisor multiplicity and
applying the sampling lemma gives
\[
 \sum_{\text{shifts}}|Z_t(C_0\Gamma\alpha)|^4
 \ll P^{2+\eps}(Q^7/L+H_*+P^2+L)
 \ll H_*P^{4+\eps-z_4(\ell)},
 \qquad z_4(\ell)=\min\{\ell,2,5-7T\}.
\]
Here $L\asymp P^\ell$ and $L^2\ll Q^7$. The normalized quadratic savings are
thus $z_2=1$ and $z_4(\ell)$.

Choose nonnegative weights $\theta_2,\theta_4,u_{i,a}$ satisfying
\[
 2\theta_2+4\theta_4\le1,\qquad
 \sum_a2a u_{i,a}\le1\quad(i\le j),\qquad
 \theta_2+\theta_4+\sum_{i,a}u_{i,a}\le1.
\]
H\"older, with missing loads estimated trivially, gives the physical saving
$\sigma(\ell)=\theta_2+\theta_4z_4(\ell)+\sum u_{i,a}b_{i,a}(\ell)$.
Substituting into \eqref{eq:unseparated} proves
\begin{equation}\label{eq:quadratic-pointwise}
 |F_j|\ll P^\eps A\,
                  P^{-\min\{1,\sigma(\ell)\}/2^{5-j}}.
\end{equation}
The cap retains the ordinary-difference diagonal. At $j=5$ it is a valid weakening
of the direct factorization. The old bound is recovered by allocating half the
mass to the quadratic second moment. Each fixed new recipe remains concave in
$\ell$, so the exact interval-endpoint verification still applies; the auxiliary
major-arc terminal is unchanged.

The current certificate uses quadratic resources at depths three, four, and five,
with the auxiliary mean at depth five for order 22. In the depth-three policy,
at its top heights,
$\ell\ge2$, $T\le3/7$, and $\phi_i\ge T-2/7$. The grouped fourth coefficient
moment then has saving $2\phi_i$, since the degree-fourteen fourth moment is
diagonal-optimal. Taking $\theta_4=u_{i,2}=1/4$ for all three coefficients gives
\[
 \sigma=(1+T)/2,\qquad
                   |F_3|\ll P^\eps A P^{-(1+T)/8}.
\]
This avoids the small sixth-moment coefficient excess in the preceding allocation.
The resulting gain is small and leaves the thirty-variable deficit positive.

\subsection{A stronger quadratic sampling estimate}

A final refinement uses the distribution of short quadratic rational heights.
Let $H=H_*$, $V=HP^2$, and let $L$ be the intrinsic height of $\alpha$ at
conductor $V$. For $B(\beta)$ the intrinsic height at conductor $P^2$, one
quadratic difference and reciprocal transference give, uniformly in the completion
parameter $t$,
\[
 |Z_t(\beta)|^2\ll P^\eps(P+P^2/B(\beta)),\qquad
 |Z_t(\beta)|^4\ll P^\eps(P^2+P^4/B(\beta)^2).
\]
The linear modulation disappears from the absolute reciprocal-sum bound, and
$B(\beta)\ll P$ by Dirichlet.

Let $N(Y)$ count $h\le H$ with $B(h\alpha)\le Y$, for $1\le Y\ll P$.
A witnessing denominator $q\le Y$ satisfies $\|hq\alpha\|\ll Y/P^2$.
Group $n=hq\ll HY$ with divisor multiplicity and apply reciprocal transference
with $J=P^2/Y$. The conductor $(HY)J=HP^2$ is exactly $V$, and hence
\[
 N(Y)\ll P^\eps\left(HY^2/P^2+HY/L+1+LY/P^2\right).
\]
Lifting the short approximant to one for $\alpha$ shows $N(Y)=0$ unless
$Y\gg L/H$. On the nonempty range the constant term is absorbed by $HY/L$;
$L^2\ll HP^2$ also absorbs the final term. Dyadic summation from
$Y\asymp\max(1,L/H)$ proves
\[
 \sum_{h\le H} B(h\alpha)^{-2}
 \ll P^\eps\left(H/P^2+\min\{H/L,H^2/L^2\}\right).
\]
Consequently
\begin{equation}\label{eq:quadratic-distribution}
 \sum_{h\le H}|Z_t(h\alpha)|^4
 \ll HP^{4+\eps}\left(P^{-2}+\min\{L^{-1},HL^{-2}\}\right).
\end{equation}
Product-of-shifts multiplicities cost only $P^\eps$, so this applies to the same
completed quadratic factors as before.

For $H=P^h$ and $L\asymp P^\ell$, the new normalized saving is
$\min\{2,\max(\ell,2\ell-h)\}$. Its upper envelope is not globally concave.
For finite certification one may instead keep the two separately valid resources
\[
 z_4^{\rm height}(\ell)=\min\{2,\ell\},\qquad
 z_4^{\rm refined}(\ell)=\min\{2,2\ell-h\}.
\]
Each is concave, so frozen H\"older recipes retain the endpoint proof. The improved
depth-five supremum can also be combined with the auxiliary second moment in the
next section. This improves the smooth 22nd moment and propagates to the higher
orders. The fully certified thirty-variable deficit nevertheless remains positive.

\section{Auxiliary means and fractional coefficient orders}

\subsection{An improved quadratic auxiliary second moment}

Let $N_i$ count the coefficient equation
$\sum_{l\le i}h_l^2(m_l^{14}-n_l^{14})=0$. The divisor argument underlying
\cite[Lemma~3.8]{VW95} gives
\[
 N_i\ll P^\eps\left(H_iM_i N_{i-1}+
                              \prod_{l<i}H_lM_l^2\right).
\]
Indeed, the zero final difference contributes the first term; for a fixed nonzero
integer $b$, the equation $h^2(m^{14}-n^{14})=b$ has $P^\eps$ solutions by divisors
and factorization. This argument needs no ordering of the $\phi_i$.

With $T_i=\sum_{l\le i}\phi_l$, the conditions
$T_i+5\phi_i\le1$ imply $N_j\ll P^\eps\widetilde H_j\widetilde M_j$.
At the quadratic terminal $j=5$, the relevant case of \cite[Lemma~3.3]{VW95}
then yields
\[
 \int_0^1|F_5|^2\dd\ll P^{1+\eps}\widetilde H_5\widetilde M_5.
\]
Combining this with the height supremum and a restricted smooth 18th moment,
using outer exponents $a=1/3$, $c=2/3$, $\gamma=1/3$, improves the smooth 22nd
moment. In general the identities are $a+c=1$, $\gamma+pa=1$, and $cv=2r$.
They give the shell inequality
\[
 \int_E|F_j||g_Q|^{2r}\dd\le
 (\sup_E|F_j|)^\gamma
 (\int_0^1|F_j|^p\dd)^a(\int_E|g_Q|^v\dd)^c.
\]

The broader source theorems contain ordering hypotheses in some cases. Those were
not silently removed: pointwise coefficient policies are order-free, whereas any
use of an ordered auxiliary theorem retains or recomputes its conditions. A full
one-block/two-block divisor recursion, with controlled excess rather than only
diagonal feasibility, is also recorded and tested in the archive.

\subsection{Real powers in the sampling lemma}

The coefficient large sieve can be extended to a real power $p\ge1$ without claiming
arbitrary-weight fractional moment estimates. Suppose its L2 sampling constant at
frequency length $M$ is $L(M)$. Let
\[
 K_M(t)=M^{-1}\left|\sum_{n=1}^M\e(nt)\right|^2,
 \qquad V_N=2K_{2N}-K_N.
\]
Applying the L2 inequality to the normalized Dirichlet kernel gives
$\sum_zK_M(x_z-t)\le L(M)$ uniformly in $t$.
The kernel $V_N$ reproduces every polynomial with frequencies in $[-N,N]$,
has $L^1$ norm at most 3, and satisfies $|V_N|\le2K_{2N}+K_N$. Jensen therefore gives
\[
 \sum_z|g(x_z)|^p\ll_p L(2N)\int_0^1|g(t)|^p\,dt.
\]
The original rational-approximation and reciprocal-sum arguments then apply unchanged.

The actual coefficient supports were checked in the source: they are the complete
smooth intervals $C(M)$, so the needed real moments follow from
$g=G_{MR}-G_M$ and the triangle inequality. An arbitrary smooth subset would require
an additional argument; no such extension is claimed.

Wooley's fractional recursion \cite[Lemma~1]{W00} and the A/B terminal schemes
of \cite[Theorem~3.4 and Section~4]{VW00} were applied in a joint finite rational
iteration to the published degree-fourteen starting moments. They give, for example,
\[
 \lambda^{(14)}_3\le\CoefficientThree,\quad
 \lambda^{(14)}_4\le\CoefficientFour,\quad
 \lambda^{(14)}_5\le\CoefficientFive.
\]
The 240 coefficient rounds and the seventh-power propagation are separately certified.
The A/B terminal uses real unweighted smooth moments only in H\"older's inequality;
all intermediate counted moments remain integral. The geometric and size hypotheses
of each auxiliary moment are checked exactly, including the distinct cases of
Theorem~3.4. No unproved optimal high-order endpoint is appended to the input table.
These are improvements of the inputs used in this calculation, without a claim
that their priority has been established in the literature.

\section{Global ordinary-moment recycling}\label{sec:recycle}

An additional improvement follows from using the ordinary 56th mean globally through
a valid counting comparison. For integers $a,s$ with $1\le a<28$ and
$u=2s-2a\ge20$, put $\gamma=1-a/28$ and $v=u/\gamma$.

\begin{proposition}
If $\Delta_v$ is admissible, then
\[
 \Delta'_{2s}=\max\{0,\gamma\Delta_v-a/28\}
\]
is admissible for the complete smooth mean.
\end{proposition}
\begin{proof}
On the complete circle, inclusion of the smooth set in the ordinary set and
orthogonality give
$\int|g_P|^{2s}\le\int|f_P|^{2a}|g_P|^u$.
This is a comparison of counts with nonnegative coefficients, not a pointwise
comparison and not a restricted-arc comparison. On the ordinary minor arcs,
\eqref{eq:f56} and H\"older give
$P^{2s-7+\gamma\Delta_v-a/28+\eps}$.
On the ordinary major arcs, $|f_P|^{2a}\le P^{2a}$ and the optimal restricted
$u$th smooth mean give $P^{2s-7}$.
For $u\ge20$, the latter follows from \cite[Theorem~1.3]{BW24}, using base order
$u-2$, height parameters 1, and cutoff 1; the central interval has measure $O(P^{-7})$.
Combining the two estimates proves the proposition.
\end{proof}

For $s=15$, $a=2$, the rule reads
\[
 \Delta_{30}\le\max\{0,(13\Delta_{28}-1)/14\}.
\]
This explains why the improved global 30th moment can have the same positive excess
as the remaining representation deficit. A positive excess is sufficient for an
admissible global mean; the representation argument requires an actual minor-arc
saving.

\subsection{The strong 32nd moment}

Globally $\int|g_P|^{32}\le\int|f_P|^2|g_P|^{30}$.
The ordinary minor-arc supremum $P^{1-1/42+\eps}$ gives a contribution
$P^{25+\Delta_{30}-1/21+\eps}$, which is smaller than $P^{25}$ because
$\Delta_{30}<1/21$. On the major arcs, the second conclusion of
\cite[Theorem~1.3]{BW24}, with base 28, order 30 and cutoff 1, gives
$\int_\M|g_P|^{30}\ll P^{23}$ without an $\eps$ loss. Therefore
\[
                    \boxed{\int_0^1|g_P|^{32}\dd\ll P^{25}.}
\]
One sufficiently small fixed smoothness exponent works for this conclusion.

Annular excesses can be negative even though complete-circle excesses are
nonnegative. The published annular rule \cite[Section~5]{BW23} permits the following
bound using the fixed certified choice $\tau_0$:
\[
 \Delta_v^*=\min_{0\le u\le v-2}(\Delta_{v-u}-u\tau_0),\qquad
 \tau_0=\frac{7-2(2.0143820)}{144}.
\]
In particular the current bound from order 30 gives
$\Delta_{32}^*\le\CurrentDeltaThirtyTwoStarExact
\approx\CurrentDeltaThirtyTwoStar$.
These annular inputs were also included in the subsequent searches.

\subsection{An optimal real moment below order 31}

There is also a useful passage from the annular estimate to a full real moment.
For every fixed $v\ge18$ and any finite set of available lower moment inputs, put
\[
 d_v=\min_{2\le w\le v}\{\Delta_w-(v-w)\tau_0\}.
\]
Then $\max(0,d_v)$ is an admissible global excess. To see this, use
\cite[Theorem~1.3]{BW24} on $\M_P(P)$, and sum
\cite[Theorem~5.2]{BW23} over the remaining dyadic annuli. If $d_v\ge0$ the
largest annular exponent is at the top height and is at most $d_v$; if $d_v<0$
the lower endpoint controls the sum. The top $R$-dependent annuli omitted from
the stated range of Theorem~5.2 are covered by its companion supremum estimate
times the lower global moment, exactly as in the proof of Theorem~5.3.

In particular,
\[
 v_0=30+\frac{\Delta_{30}}{\tau_0}
       =\CurrentRealZeroExact\approx\CurrentRealZero
\]
has admissible excess zero. The endpoint permits an $\eps$ loss. At every fixed
$v>v_0$ the minor arcs have a strict saving and the major estimate has no
$\eps$ loss. This proves, for example,
\[
                     \boxed{\int_0^1|g_P|^{31}\dd\ll P^{24}.}
\]
This is a pure smooth moment. Mixed ordinary/smooth moments require their own
argument, and no minor-arc saving at order 30 follows from this statement.

\section{What the finite certificate establishes}

The current certificate uses \CertificateRounds{} finite rounds, with rational
parameters, rational H\"older weights, and upward rounding of intermediate moment
values. It checks the full height cover, each moment load, every size condition,
the separate auxiliary major terminal, and every efficient-differencing recursion.
The root independently reproduced the stored outputs of the successive exact
certificates. The analytic hypotheses were reviewed independently of those numerical
reruns.

The following are parameter summaries, rounded here only for readability. The exact
values and all height recipes are in the JSON certificate.
\begin{center}\small
\begin{tabular}{r l l l}
\toprule
Order & Inner half-moment path & Size exponents $\phi_i$ & Terminal type\\
\midrule
\ParameterRows
\bottomrule
\end{tabular}
\end{center}

Two bookkeeping issues were explicitly corrected. An unused appended
degree-fourteen endpoint claiming $\Delta_{82}=0$ was replaced by the valid trivial
extension from order 80. The selected earlier certificates used coefficient half-orders
at most 25, so that endpoint did not affect their conclusions. Also, the legacy field
name \texttt{two\_ordinary} had denoted the seven-ordinary expression
$(6\Delta_{26}+\Delta_{28}-1)/8$ in some earlier files. New certificates use the
correct labels. Neither correction changes the established main deficit.

No floating-point search is treated as proof of optimality. Exact dual checks in the
archive concern their explicitly frozen H\"older families only. They do not exclude
stronger estimates, a different representation, or another proof of $G(7)\le30$.

\section{Real mixed products and the finite comparison}
\label{sec:poisson}

The correlation inputs can be extended to real mixed powers by a positive
Poisson majorant. Let $H_i$ be polynomials of degrees $N_i$, let $a_i>0$, and put
\[
 U(t)=\prod_i|H_i(\e(t))|^{a_i},\qquad D=\sum_i a_iN_i>0.
\]
If $\mathcal P_r$ is the circle Poisson kernel, subharmonicity gives
$U\le\mathcal P_{1/R}*U_R$ on a disk of radius $R>1$.
Applying the same inequality to the conjugate reversed polynomials gives
$U_R\le R^D\mathcal P_{1/R}*U$. The semigroup identity and $R=\exp(1/D)$ prove
\[
                 U\le e\,\mathcal P_{\exp(-2/D)}*U.
\]
This argument needs no nonnegative Fourier expansion of a fractional power.

For $L_P=\sum_{P<x\le2P}\e(\alpha x^7)$, pair spacing gives
$|x^7-y^7|\ge7P^6|x-y|$. The Poisson convolution of $|L_P|^2$ is consequently
$O(P+DP^{-5})$. The Robert--Sargos spacing estimate gives the corresponding
fourth convolution bound $O(P^\eps(P^2+DP^{-3}))$. Fubini therefore proves
\begin{align*}
 \int|L_P|^2U&\ll(P+DP^{-5})\int U,\\
 \int|L_P|^4U&\ll P^\eps(P^2+DP^{-3})\int U.
\end{align*}
For a product of fixed real powers of ordinary and smooth sums at scale $Q$,
one has $D\ll Q^7$. Both comparisons use complete-circle moments on the right.

These inputs were included in a finite comparison with 3,935 columns:
ordinary moments, current pure and mixed smooth moments, their valid correlations,
pointwise estimates, and direct ordinary-minor-mean recycling. The variable counts
were allowed to be real, which enlarges the options beyond integer allocations.
Exact dual witnesses give the following uniform positive lower bounds for the
exponent deficit when $Q=P^\theta$:
\begin{center}
\begin{tabular}{lll}
\toprule
Allowed counts totaling 30 & Range of $\theta$ & Deficit at least\\
\midrule
Nonnegative real counts & $0\le\theta\le6/7$ & $\ConeGapShort$\\
At least 9 long ordinary & $0\le\theta\le6/7$ & $\ConeGapShortNine$\\
Nonnegative real counts & $0\le\theta\le1$ & $\ConeGapAll$\\
At least 9 long ordinary & $0\le\theta\le1$ & $\ConeGapAllNine$\\
\bottomrule
\end{tabular}
\end{center}
The input costs are affine between $\theta=0,5/7,6/7,1$.
Feasible exact endpoint duals interpolate over these intervals. An independent
standard-library verifier checks every rational column inequality, with no
floating optimizer. This excludes a power saving within the listed finite family;
it is not an impossibility theorem for other analytic estimates.

The final comparison also allows every H\"older factor to choose its own short
scale $Q_i=P^{\theta_i}$. For a column with counts $(a,b,c)$, subtract its main-scale
contribution $a+\theta_i(b+c)$ from its moment exponent. The objective is then
$7+\sum w_i\,\mathrm{cost}_i$, with total variable count 30 and H\"older mass at
most one. The union of the four endpoint scales has $4\times3935=15740$ columns.
One exact dual works simultaneously at all four endpoints, hence at every
independently selected $\theta_i\in[0,1]$ by affine interpolation. The same lower
bounds $\ConeGapAll$ and $\ConeGapAllNine$ remain, respectively without a long-count
restriction and with at least nine long variables. This is an enlarged algebraic
comparison: additional common-arc and representation requirements can only restrict
its feasible choices. Its positive result excludes this extension of the specified
input family as well.

\section{A second route: four long variables correlated with shorter sums}\label{sec:weighted}

Write $L_P(\alpha)=\sum_{P<x\le2P}\e(\alpha x^7)$ and $Q=P^\theta$.
The Robert--Sargos spacing estimate \cite[Theorem~2]{RS06} gives
\[
 N(H)\ll P^\eps(P^2+P^{-3}H)
\]
for near coincidences of two sums of two long seventh powers. For even $v$, the
uniform coefficient bound from a smooth mean yields the following estimate;
Section~\ref{sec:poisson} supplies its real-power extension:
\[
 \int_0^1|L_P|^4|g_Q|^v\dd
       \ll P^\eps(P^2+P^{-3}Q^7)Q^{v-7+\Delta_v}.
\]
The excess $\Delta_v$ in the second term is the unresolved correlation cost.

\subsection{Two half-excess hypotheses}

At $\theta=747/1000$, assume, for $v=20,22$,
\begin{equation}\label{eq:Hv}
 \int_0^1|L_P|^4|g_Q|^v\dd\ll P^\eps
       \left(P^2Q^{v-7+\Delta_v}+P^{-3}Q^{v+\Delta_v/2}\right),
\end{equation}
using the published $\Delta_{20}=.6055676$, $\Delta_{22}=.4242973$.
Then H\"older with weights $9/52$, $31/52$, $12/52$ on the ordinary 56th mean and
the two mixed means controls $L_P^{13}g_Q^{17}$. The saving below its main scale
$P^6Q^{17}$ is exactly
\[
 \frac{54209403}{32500000000}=.001667981630769\ldots>0.
\]
The complete conditional circle-method implication was checked, including the
ordinary interval minor arcs, pruning, positive arithmetic factor, and the interior
real point. Here the minor domain is outside the union of the ordinary major sets
for $f_P$ and $f_{2P}$, or outside a fixed-constant enlargement $\M_P(CP)$.
Then $L_P=f_{2P}-f_P$ inherits the required 56th minor moment.
It uses 13 long and 17 short variables, hence 30 in total.

For the real factor choose $P=(n/[13(3/2)^7])^{1/7}$, placing the 13 long real
coordinates strictly inside $(P,2P)$. The short contribution is $o(P^7)$.
At narrow arcs the arithmetic singular series still has 30 variables, while the
limiting real integral has 13 long variables. Confusing these two counts would be an
error. Ordinary major-arc pruning is available with 13 long variables alone.

\subsection{One explicit 22nd mixed-moment hypothesis}

A simpler sufficient input, at $Q=P^{37/50}$, is
\begin{equation}\label{eq:H22}
 \int_0^1|L_P|^4|g_Q|^{22}\dd\ll P^\eps
       \left(P^2Q^{15.4242973}+P^{-3}Q^{22.19}\right).
\end{equation}
H\"older weights $83/616$ for the ordinary 56th mean, $19/22$ for
\eqref{eq:H22}, and $1/616$ for measure give 11 long and 19 short variables.
The exact saving is $1501/770000=.001949350649\ldots$.
Again the major arcs and positivity complete the conditional implication.

Neither \eqref{eq:Hv} nor \eqref{eq:H22} is proved on the complete remaining minor
region. The following partial regions are proved:
\begin{itemize}
\item The two half-excess estimates on $\M_Q(Q^{1.745})$, with a slightly larger
      outer buffer for the local reproducing kernel.
\item The single22nd-moment target on $\M_Q(Q^{1.56})$.
\item For that single target, also the Dirichlet-scale arcs
      $q\le Q^{.095}$, $|q\alpha-a|\le Q^{-7/2}$.
\item The long-variable slice $x_1+x_2=x_3+x_4$, whose nontrivial near-coincidence
      count is $\ll P^{-4+\eps}Q^7$.
\end{itemize}
The low-denominator extension uses Bourgain's translated local20th moment
\cite[Section~6]{Bou16}, valid also for the smooth indicator, and the local
fourth-moment operator obtained from Robert--Sargos. These proved pieces do not
cover the unresolved deep region.

\newpage
\section{Two sufficient pure smooth-moment improvements}\label{sec:pure-targets}

The following are separate sufficient hypotheses. Neither is proved in this
investigation, and neither assumes the other. They concern only the pure smooth
sum, not an additional mixed ordinary/smooth estimate:
\begin{align}
 \int_0^1|g_P|^{22}\dd&\ll P^{15.31+\eps},\label{eq:pure22-target}\\
 \int_0^1|g_P|^{20}\dd&\ll P^{13.515+\eps}.\label{eq:pure20-target}
\end{align}
As usual each is required for every positive $\eps$ and sufficiently small fixed
smoothness. The certified unconditional excesses are respectively
$\Delta_{22}=\CurrentDeltaTwentyTwo$ and $\Delta_{20}=\CurrentDeltaTwenty$;
the hypotheses would reduce
them to $.31$ or $.515$.

For each hypothesis separately, an exact 220-round calculation starts with the
published moment table and changes precisely one entry. It then uses the frozen
grouped-cubic policies, the existing auxiliary estimates, convexity, and ordinary
moment recycling. The conditional certificates preserve the original table and
the single override explicitly. An independently written verifier reproduced
both calculations. Their conclusions are
\begin{center}
\begin{tabular}{lll}
\toprule
Unproved input & Resulting admissible $\Delta_{28}$ & Four-ordinary deficit\\
\midrule
$\Delta_{22}\le31/100$ & $\le.075329303761$ &
 $-20719051107/14000000000000$\\
$\Delta_{20}\le103/200$ & $\le.074844741085$ &
 $\le-.001929883278214$\\
\bottomrule
\end{tabular}
\end{center}
The second row is rounded in the safe direction. Both excesses are strictly below
$1/13$, so \eqref{eq:gap} gives a minor-arc power saving. The established major-arc
asymptotic and positivity then give $G(7)\le30$ under either hypothesis. The
conditional implications use the earlier grouped certificate; the later quadratic
fourth-moment refinement is not needed for them.

The distinction is essential: these computations identify sufficient analytic
targets. They do not prove either target and do not improve the unconditional
record.

\section{Precise ordinary minor-arc alternatives}\label{sec:ordinary}

The following numerical implications use the frozen enhanced-recycling certificate
with $\Delta_{26}=.176210893849$ and $\Delta_{28}=.093725806221$; any later smaller
admissible values preserve them.

\begin{proposition}[Conditional 52nd-moment route]
If
\[
 \int_{\m_P(P^{17/10})}|f_P|^{52}\dd\ll P^{44+\eps},
\]
then the certified smooth moments imply $G(7)\le30$.
\end{proposition}
\begin{proof}
Use six ordinary and 24 smooth variables. On the indicated deep set H\"older has
weights $3/26$ and $23/26$, and effective smooth order $624/23$.
The saving is
\[
 \frac{3-10\Delta_{26}-13\Delta_{28}}{26}
       =\frac{19455580637}{26000000000000}>0.
\]
On $\M_P(P^{17/10})\setminus\M_P(P)$, the ordinary supremum and the restricted
24th smooth mean save $(1-(17/5)\Delta_{24})/7=.005051275973\ldots>0$.
Optimal smooth major-arc pruning and the usual positive 30-variable singular factors
complete the argument.
\end{proof}

Similarly, the hypothesis
\[
 \int_{\m_P(P^{8/5})}|f_P|^{56}\dd\ll P^{48-4/25+\eps}
\]
suffices using seven ordinary and 23 smooth variables. The deep saving is
$1801766137/1600000000000=.001126103835625$, and the intermediate saving is
$.005027383396\ldots$. More generally this allocation needs an improvement
$\beta>.150991169315$ in the exponent 48.

These are explicit sufficient targets, not newly proved ordinary estimates. The
known 56th minor mean and Hua 16th mean interpolate only to a 52nd exponent of 44.4.
The additional geometric sixth-moment estimate examined in this run does not
improve that particular interpolation.

\section{Other investigations and their exact limits}

\textbf{Additional quadratic resources.}
The height-distribution proof also gives real sampled moments between orders two
and four. A third-moment check at the final frozen parameters gave no gain at the
controlling endpoint. At the intrinsic top height, the new fourth-moment resource
has saving two and dominates the third-moment alternative there. Within the same
separate-coefficient allocation, higher quadratic moments and a quadratic supremum
are dominated as well: every coefficient saving is at most $h=5-7T$, and the
additional fourth-moment saving exactly compensates the H\"older mass needed to
replace those alternatives. This conclusion concerns these resources and allocations.

\textbf{Intermediate fractional differencing.}
For nonnegative integer $u$, real $h\ge1$, $0\le\theta\le1/2$, and
$r=\theta u+(1-\theta)h$, the valid shell rule is
\[
 \int_E|F_j||g|^{2r}\le(\sup_E|F_j|)^{1-2\theta}
       \left(\int_0^1|F_j|^2|g|^{2u}\right)^\theta
       \left(\int_E|g|^{2h}\right)^{1-\theta}.
\]
The counted mixed square remains global. The height supremum, restricted smooth
mean and ordinary next-step recursion can all be retained. Tested policies,
including annular inputs and multiple auxiliary moments, did not close the gap.

\textbf{The unrestricted linear terminal.}
At differencing depth six, the centered phase is exactly
$2^6\cdot7!\,(\prod h_i)z$. A positive Fourier majorant for the Dirichlet kernel,
followed by a divisor count for each nonzero Fourier frequency, proves
\[
 \int_0^1|F_6||g_Q|^{2r}\dd
 \ll P^\eps\widetilde M\{Q^{2r}+\widetilde H\int_0^1|g_Q|^{2r}\dd\}
\]
for integer $r$, without the earlier height restriction. Its wider validity was
checked against the actual Cartesian summation ranges. The tested larger parameter
region gave no improvement; a further convexity calculation explains why the
associated counted-square terminal is dominated by the available pointwise-height
terminal in the specified interpolation family.

\textbf{Denominator arithmetic.}
Retaining $\nu_3(q)=q^{-1}\sum_{d\mid q}d/\rho_3(d)$, where
$\rho_3(d)=\prod_{p^a\parallel d}p^{\lceil a/3\rceil}$, improves the cubic reciprocal
weight for many denominators. A valid adaptation of the weighted pruning lemma gives
\[
 T^\eps X^{2t-7}+X^{-7}T^{t+1+\eps}.
\]
Here $t\ge3$, the rational box satisfies $TY\le X^7/4$, and the beta-weight
exponent is integrable; the exponent-1 endpoint requires a logarithmic treatment.
Its completion term becomes too large at $T=X^{7/2}$. It cannot be suppressed
as an $\eps$ loss. The immediate beta decay also requires an additional cost before
an integrable pruning weight is available.

\textbf{Long translations.}
Extending the averaging range of translated moment-curve sums does not produce
a free power saving. An integer unipotent change of the coefficient variables
identifies each translated block's weighted integral with the original one;
the number of blocks cancels the longer averaging normalization. A genuine
collective maximal or restricted-domain estimate is still required.

\textbf{Recent small-cap and inhomogeneous results.}
The recent slab estimates were checked for the actual coefficient domains needed
here \cite{JO26,OY25}.
A bound for a slab is not a bound for an arbitrary union of tilted strips. Fixed
inhomogeneity results also do not supply uniform control for the growing family of
inhomogeneities in the transfer. Summing their individual bounds loses more than
the saving they provide. Announced work without an available complete proof was
excluded from certified inputs.

\textbf{Geometric low moments.}
The literature branch derived an ordinary sixth-moment upper exponent
$12/5+27/(10\sqrt7)=3.4205040771\ldots$ from published geometric inputs
\cite{Sal05,Sal23,Ver25}, with
the singular-section branches checked separately. It improves the elementary
sixth-moment interpolation, but its interpolation with 56th exponent 48 gives
52nd exponent 44.4336\ldots, weaker than the existing 44.4 line. The complete
derivation and its source audit are retained as a secondary result, outside the
core record calculation.

\section{Conclusion and the next mathematical decision}

The established record remains $G(7)\le\BoundGSeven$. The work has produced
reproducible auxiliary progress: a sharpened coefficient reciprocal estimate,
height-coupled integration, direct cubic-product grouping, a retained quadratic
fourth moment, a lower auxiliary moment, real-power coefficient sampling,
improved degree-fourteen inputs, and stronger global smooth-moment recycling.
Each must be distinguished from the still missing representation saving.

The most concrete remaining choices are an analytic improvement that pushes
$\Delta_{28}$ below $1/13$, either sufficient pure-moment improvement in
Section~\ref{sec:pure-targets}, one of the ordinary deep-minor moment targets in
Section~\ref{sec:ordinary}, or the four-long/short correlation estimates in
Section~\ref{sec:weighted}. Additional
optimization within the families already tested is a narrower prospect: the
recorded numerical barriers concern those families only, but the repeated small
gains do not justify announcing a record proof.

\appendix
\section{Reproduction and audit map}

The archive contains the current certificate and its generating script, the
degree-fourteen fractional certificate, exact arithmetic for the conditional
targets, and independent source-level audits. The exact scripts use Python's
\texttt{Fraction} arithmetic. Search scripts additionally use numerical optimization;
their outputs are explicitly marked as exploratory.

Run \texttt{python3 verify\_results.py} from the extracted archive to replay the
degree-fourteen updates, the current seventh-power certificate, both explicitly
conditional pure-moment implications, and the finite-cone witness. This command
needs only the Python standard library. The selected seventh-power certificate is
\CertificateName.

\begin{minipage}{\linewidth}
The degree-fourteen certificate is\par
\mbox{\small\path{weighted_correlation/degree14_joint_curve_certificate.json}}.
The ordinary-target arithmetic is in\par
\mbox{\small\path{root/ordinary_minor_target_certificate.py}}.
\end{minipage}

The main analytic notes are under \path{root/}, with independent reviews under
\path{weighted_correlation/} and \path{literature_routes/}. They retain the hypotheses,
scope corrections, and negative results needed to avoid repeating invalid transfers.
The archive manifest records file sizes and SHA256 digests, and the archive was
checked by reopening it and verifying its contents.

\begin{thebibliography}{20}\small
\bibitem{W16} T. D. Wooley, \emph{On Waring's problem for intermediate powers},
Acta Arith. 176 (2016), 241--247.
\url{https://arxiv.org/abs/1602.03221}.
\bibitem{VW95} R. C. Vaughan and T. D. Wooley, \emph{Further improvements in
Waring's problem}, Acta Math. 174 (1995), 147--240.
\url{https://www.math.purdue.edu/~twooley/publ/1995%20fiwp.pdf}.
\bibitem{VW00} R. C. Vaughan and T. D. Wooley, \emph{Further improvements in
Waring's problem, IV: Higher powers}, Acta Arith. 94 (2000), 203--285.
\url{https://www.math.purdue.edu/~twooley/publ/2000%20fiwp4.pdf}.
\bibitem{W00} T. D. Wooley, \emph{Quasi-diagonal behaviour and smooth Weyl sums},
Monatsh. Math. 130 (2000), 161--170.
\url{https://www.math.purdue.edu/~twooley/publ/2000SWSQD.pdf}.
\bibitem{BW23} J. Br\"udern and T. D. Wooley, \emph{On Waring's problem for larger
powers}, J. reine angew. Math. 805 (2023), 115--142.
\url{https://www.math.purdue.edu/~twooley/publ/20230915wpk.pdf}.
\bibitem{BW24} J. Br\"udern and T. D. Wooley, \emph{Estimates for smooth Weyl sums
on major arcs}, 2024.
\url{https://arxiv.org/abs/2405.18608}.
\bibitem{RS06} O. Robert and P. Sargos, \emph{Three-dimensional exponential sums
with monomials}, J. reine angew. Math. 591 (2006), 1--20.
\url{https://perso.univ-st-etienne.fr/rool6510/robert-2006-crelle.pdf}.
\bibitem{Bou16} J. Bourgain, \emph{On the Vinogradov mean value}, 2016,
Theorem 10 and Section 6. \url{https://arxiv.org/abs/1601.08173}.
\bibitem{Sal05} P. Salberger, \emph{Counting rational points on hypersurfaces
of low dimension}, Ann. Sci. ENS 38 (2005), 93--115.
\url{https://www.numdam.org/item/ASENS_2005_4_38_1_93_0.pdf}.
\bibitem{Sal23} P. Salberger, \emph{Counting rational points on projective
varieties}, Proc. London Math. Soc. 126 (2023), 1092--1133.
\url{https://doi.org/10.1112/plms.12508}.
\bibitem{Ver25} M. Verzobio, \emph{Counting rational points on smooth
hypersurfaces with high degree}, IMRN (2025), rnaf249.
\url{https://arxiv.org/abs/2503.19451}.
\bibitem{JO26} S. Jeong and C. Oh, \emph{Amplitude dependent wave envelope
type estimates for moment curves and applications}, 2026.
\url{https://arxiv.org/abs/2611.26581}.
\bibitem{OY25} C. Oh and K. Yeon, \emph{An extended Vinogradov's mean value
theorem}, 2025. \url{https://arxiv.org/abs/2506.01751}.
\end{thebibliography}
\end{document}
